\section*{Chapter \ref{ch:mx:market-quality-and-regulation} --- Market Quality and Regulation}

\begin{solution}{exo:mx:market-quality-and-regulation:1}
$(3.18-1.58)-(2.27-1.56)=1.60-0.71=0.89$ basis points.
\end{solution}

\begin{solution}{exo:mx:market-quality-and-regulation:2}
$900/250=3.6$ messages per trade; with the rule's effect, about 2.4.
\end{solution}

\begin{solution}{exo:mx:market-quality-and-regulation:3}
Trades alternate between bid and ask, so short-interval returns are negatively autocorrelated and their variances add up to more than the long
interval's. The rule lowered the ratio by 0.14 (0.13 estimated): wider spreads, more bounce.
\end{solution}

\begin{solution}{exo:mx:market-quality-and-regulation:4}
A stock's days share its unobserved characteristics and its errors are correlated over time; the rule is assigned by stock, so the effective number
of independent observations is the number of stocks, not of stock-days.
\end{solution}

\begin{solution}{exo:mx:market-quality-and-regulation:5}
It targeted 9\% of the previous minute's volume, and the intermediaries' \omterm{def:mx:market-quality-and-regulation:hotpotato}{hot-potato trading} inflated that volume without absorbing any of the selling:
the algorithm sold faster into a market that was not buying.
\end{solution}

\begin{solution}{exo:mx:market-quality-and-regulation:6}
Order-book data show only the orders that changed the book: the losers of a race (failed trades and cancels) leave no trace. Exchange message data,
with every message and its outcome, do.
\end{solution}

\begin{solution}{exo:mx:market-quality-and-regulation:7}
1.11 basis points (standard error 0.06) against a truth of 0.94: the least active stocks suffer more from the market-wide change than the controls,
so the control group understates what the treated stocks' spreads would have done without the rule; parallel trends fail.
\end{solution}

\begin{solution}{exo:mx:market-quality-and-regulation:8}
The 1.6 basis points include the market-wide change that hit every stock (0.71 on the controls); the rule's own effect, by \omterm{def:mx:tick-size-queues-and-priority:did}{difference-in-differences},
is 0.89, close to the true 0.87.
\end{solution}

\begin{solution}{pb:mx:market-quality-and-regulation:1}
\textbf{1.} Costs, depth and price informativeness: quoted, effective and realised spreads, impact, depth, resilience, Amihud, the variance ratio.
\textbf{2.} Effective = realised + price impact.
\textbf{3.} See the definitions.
\textbf{4.} How far prices are from a random walk over short intervals: below one with bid-ask bounce.
\textbf{5.} \omterm{def:mx:market-quality-and-regulation:hft}{Algorithmic trading} narrowed spreads, reduced adverse selection and made quotes more informative for large stocks; identified by the NYSE's
2003 automation of quote dissemination as an instrument.
\textbf{6.} High-frequency traders trade in the direction of permanent price changes and against transitory errors, mostly with their aggressive orders.
\textbf{7.} 75,000 E-mini contracts at 9\% of volume without regard to price or time; high-frequency traders then traded over 27,000 contracts among
themselves in 14 seconds, about 49\% of volume, buying about 200 net.
\textbf{8.} Latency-arbitrage races (about one a minute per FTSE~100 symbol, modal duration 5--10 microseconds, about 20\% of volume), seen in message data
that record failed attempts.
\textbf{9.} 20 stocks, 16 ten-minute days; from day 8 a market-wide rise in volatility and fall in quoting for all, and an order-to-trade cap on every
other stock.
\textbf{10.} By running each treated stock-day after the event also without the rule, with the same order flow.
\textbf{11.} $y_{id}=\alpha_i+\gamma_d+\beta\,\text{treated}_i\times\text{post}_d+\varepsilon_{id}$, errors clustered by stock.
\textbf{12.} \emph{Named result}: \omterm{def:mx:tick-size-queues-and-priority:did}{difference-in-differences} 0.89 basis points (standard error 0.09) against a before-and-after 1.60 (0.08) that ignores
the control group; the truth is 0.87 (0.03).
\textbf{13.} Coefficients within 0.07 of zero before the event: the groups moved together, supporting parallel trends.
\textbf{14.} It cut messages but widened spreads (0.9 quoted, 1.3 effective), reduced depth and increased bounce: worse by every trading-cost
measure.
\textbf{15.} \omterm{def:mx:tick-size-queues-and-priority:did}{Difference-in-differences} overstates: 1.11 against 0.94.
\textbf{16.} The \omterm{def:mx:market-quality-and-regulation:otr}{order-to-trade ratio} (down about a quarter), the rule's target.
\textbf{17.} A comparable control group (or staggered adoption), pre-period data to check trends, and measures from message-level data.
\textbf{18.} \omterm{def:mx:market-quality-and-regulation:hotpotato}{Hot-potato trading} inflates it without anyone taking risk, as in the flash crash.
\textbf{19.} Venues must calculate each member's ratio of unexecuted orders to transactions for every instrument.
\textbf{20.} A measure of what would have happened anyway.
\end{solution}

\begin{solution}{iq:mx:market-quality-and-regulation:1}
Compare with stocks the change did not cover over the same days, check that they moved together before, and control for volatility and volume; a
\omterm{def:mx:tick-size-queues-and-priority:did}{difference-in-differences} with an event study.
\iqlookfor{A control group; pre-trends.}
\end{solution}

\begin{solution}{iq:mx:market-quality-and-regulation:2}
The treated group's change minus the control group's; it assumes both would have changed alike without the treatment (parallel trends), tested by
the event study's pre-period coefficients and by placebo dates.
\iqlookfor{Parallel trends; placebo tests.}
\end{solution}

\begin{solution}{iq:mx:market-quality-and-regulation:3}
In normal times, mostly yes: narrower spreads, more informative prices (Hendershott, Jones and Menkveld; Brogaard, Hendershott and Riordan). It does
not settle behaviour in stress (the flash crash), the cost of latency races (Aquilina, Budish and O'Neill), or who bears adverse selection.
\iqlookfor{Balance; stress and races.}
\end{solution}

\begin{solution}{iq:mx:market-quality-and-regulation:4}
A participation or volume-following algorithm needs a price limit and a check that volume is real liquidity; monitor the book's depth, not just volume,
and slow down when prices move fast.
\iqlookfor{Price awareness; volume versus depth.}
\end{solution}

\begin{solution}{iq:mx:market-quality-and-regulation:5}
Count messages and trades per instrument and account in the order gateway, compute the venue's ratio on its own formula intraday, throttle quoting
strategies as they approach the limit, and alert before breaches.
\iqlookfor{Real-time counting; throttles; the venue's definition.}
\end{solution}

\begin{solution}{iq:mx:market-quality-and-regulation:6}
From message data with every order, cancel and failed attempt and exchange timestamps: group messages that hit the same \omterm{def:mx:the-limit-order-book:tick}{price level} within
microseconds, with at least one winner and one loser; the venue must provide failed messages, not just the book.
\iqlookfor{Message-level data; defining a race.}
\end{solution}
